% ═══════════════════════════════════════════════════════════════════════════
%  فصل ۴ — ایران باستان: هخامنشیان و ساسانیان
% ═══════════════════════════════════════════════════════════════════════════
\chapter{ایران باستان: هخامنشیان و ساسانیان}
\label{ch:ancient-iran}

\begin{tcolorbox}[mybox, title={چکیده‌ی فصل}]
این فصل به بررسی نخستین نمونه‌های مستند نجات‌طلبی و دعوت از نیروی خارجی در تاریخ ایران باستان می‌پردازد. از دعوت‌های احتمالی ساتراپ‌ها و اقوام ناراضی از اسکندر مقدونی در دوره‌ی هخامنشی تا روابط پیچیده‌ی مرزنشینان ساسانی با رومیان و اعراب، الگوهای تکرارشونده‌ی نجات‌طلبی در بستر ایران باستان واکاوی می‌شود. تحلیل بر مبنای مدل شش‌وجهی فصل ۳ و با تکیه بر منابع یونانی، لاتین، فارسی میانه و عربی صورت می‌گیرد.
\end{tcolorbox}

% ─────────────────────────────────────────────────────────────────────────
\section{مقدمه: ساختار سیاسی و آسیب‌پذیری امپراتوری‌های باستانی}
\label{sec:ch04-intro}
% ─────────────────────────────────────────────────────────────────────────

امپراتوری‌های باستانی به‌دلیل وسعت جغرافیایی، تنوع قومی-فرهنگی و ضعف سازوکارهای ارتباطی، ذاتاً مستعد شکاف‌های داخلی بودند. \thinker{جوزف ویسهوفر} (\lr{Josef Wiesehöfer}) تأکید می‌کند که «شاهنشاهی هخامنشی، بر خلاف تصور یونانیان، یک بلوک یکپارچه نبود بلکه مجموعه‌ای از واحدهای نیمه‌خودمختار بود که وفاداری‌شان به مرکز همواره در نوسان بود.»%
\footnote{\lr{Wiesehöfer, Josef. \textit{Ancient Persia: From 550 BC to 650 AD}. Trans. Azizeh Azodi. I.B. Tauris, 2001, p.~68.}}

این ساختار، زمینه‌ساز دو نوع نجات‌طلبی در ایران باستان بود:
\begin{enumerate}[nosep, label=\textbf{\arabic*.}]
  \item \textbf{نجات‌طلبی گریز از مرکز:} ساتراپ‌ها، اقوام حاشیه‌ای یا نخبگان ناراضی که از دشمن خارجی علیه شاهنشاه یاری می‌طلبیدند.
  \item \textbf{نجات‌طلبی معکوس:} شاهنشاه یا یکی از مدعیان تاج‌وتخت که از نیروی خارجی برای سرکوب شورش داخلی یا تحکیم قدرت یاری می‌خواست.
\end{enumerate}

% ─────────────────────────────────────────────────────────────────────────
\section{دوره‌ی هخامنشی (۵۵۰–۳۳۰ ق.م)}
\label{sec:ch04-achaemenid}
% ─────────────────────────────────────────────────────────────────────────

\subsection{تأسیس شاهنشاهی: کورش و الگوی «آزادسازی»}
\label{subsec:ch04-cyrus}

خود تأسیس شاهنشاهی هخامنشی نمونه‌ای بارز از نجات‌طلبی معکوس است. بر اساس \concept{استوانه‌ی کورش} و روایت‌های بابلی، \thinker{کورش بزرگ} از سوی بخشی از کاهنان و نخبگان بابلی — که از سیاست‌های \enemy{نبونید} (آخرین پادشاه بابل) ناراضی بودند — به‌عنوان «آزادکننده» خوانده شد.%
\footnote{\lr{Kuhrt, Amélie. "The Cyrus Cylinder and Achaemenid Imperial Policy." \textit{Journal for the Study of the Old Testament} 25 (1983): 83--97.}}

\begin{histQuote}{استوانه‌ی کورش — ترجمه‌ی آزاد}
«مردوک [خدای بزرگ بابل] ... در میان همه‌ی سرزمین‌ها جست‌وجو کرد و شاهی دادگر را یافت که قلبش از او خشنود بود: کورش، شاه انشان. او را به نام خواند تا بر تمامی جهان فرمانروایی کند ... بدون نبرد و جنگ، او [مردوک] اجازه داد کورش وارد بابل شود.»%
\footnote{\lr{Finkel, Irving. \textit{The Cyrus Cylinder}. I.B. Tauris, 2013, lines 12--17.}}
\end{histQuote}

\begin{tcolorbox}[casebox, title={تحلیل شش‌وجهی: ورود کورش به بابل (۵۳۹ ق.م)}]
\begin{itemize}[nosep, rightmargin=1em]
  \item \textbf{ساختاری:} بحران مشروعیت نبونید به‌دلیل بی‌اعتنایی به آیین‌های مردوک و اقامت طولانی در تیما.%
  \footnote{\lr{Beaulieu, Paul-Alain. \textit{The Reign of Nabonidus, King of Babylon, 556--539 B.C.} Yale UP, 1989.}}
  \item \textbf{حقوقی-هنجاری:} مفهوم مدرن حقوق بین‌الملل وجود نداشت؛ هنجار غالب «حق الهی» بود. کاهنان بابلی با ادعای رضایت مردوک، مشروعیت ورود کورش را تأمین کردند.
  \item \textbf{اقتصادی:} بابل ثروتمندترین شهر جهان باستان بود. کورش با حفظ ساختار اقتصادی بابل و عدم غارت، هزینه‌ی اقتصادی مداخله را به حداقل رساند.
  \item \textbf{نظامی:} عملاً بدون نبرد جدی. ارتش نبونید در اُپیس شکست خورد و بابل بدون محاصره گشوده شد.%
  \footnote{\lr{Briant, Pierre. \textit{From Cyrus to Alexander: A History of the Persian Empire}. Trans. P.T. Daniels. Eisenbrauns, 2002, pp.~41--44.}}
  \item \textbf{روان‌شناختی:} نارضایتی کاهنان و بخشی از جمعیت از نبونید، زمینه‌ی پذیرش «ناجی خارجی» را فراهم کرده بود.
  \item \textbf{ژئوپولیتیک:} فتح بابل، کورش را به فرمانروای خاور نزدیک تبدیل کرد — منفعتی عظیم.
  \item \textbf{نتیجه (طیف):} \textbf{A–B}. بابل تا حد زیادی خودمختاری فرهنگی و اداری خود را حفظ کرد، اما حاکمیت سیاسی مستقل خود را برای همیشه از دست داد.
\end{itemize}
\end{tcolorbox}

\subsection{شورش‌های ساتراپی و دعوت از یونانیان}
\label{subsec:ch04-satrapal-revolts}

یکی از مهم‌ترین نمونه‌های نجات‌طلبی در دوره‌ی هخامنشی، \concept{شورش بزرگ ساتراپ‌ها} (\lr{Great Satraps' Revolt}) در دهه‌ی ۳۶۰ ق.م است. چندین ساتراپ آسیای صغیر — از جمله \enemy{آریوبرزانس} و \enemy{دَتَمِس} — علیه \histref{اردشیر دوم}{۴۰۴–۳۵۸ ق.م} شوریدند و از \concept{اسپارت} و \concept{مصر} درخواست کمک نظامی کردند.%
\footnote{\lr{Weiskopf, Michael. \textit{The So-Called "Great Satraps' Revolt", 366--360 B.C.} Franz Steiner Verlag, 1989.}}

\thinker{دیودوروس سیکولوس} (\lr{Diodorus Siculus}) گزارش می‌دهد که ساتراپ‌های شورشی «از فرعون مصر و از لاکدایمونیان [اسپارتی‌ها] و از سایر یونانیان خواستند که متحدشان شوند و وعده‌ی سهمی از ثروت شاهنشاهی را دادند.»%
\footnote{\lr{Diodorus Siculus. \textit{Bibliotheca Historica}. XV.90--92. Trans. C.H. Oldfather. Loeb Classical Library, 1954.}}

\begin{tcolorbox}[enemybox, title={الگوی شکست: شورش ساتراپ‌ها}]
شورش ساتراپ‌ها در نهایت شکست خورد و دلایل آن آموزنده است:
\begin{enumerate}[nosep, label=\textbf{\arabic*.}]
  \item \textbf{عدم انسجام داخلی:} ساتراپ‌ها اهداف متفاوتی داشتند — برخی استقلال می‌خواستند و برخی صرفاً امتیازات بیشتر از شاهنشاه.
  \item \textbf{بی‌اعتمادی متقابل:} یونانیان و مصریان هر یک منافع خود را دنبال می‌کردند و تعهد واقعی به ساتراپ‌ها نداشتند.
  \item \textbf{خیانت درونی:} برخی ساتراپ‌ها توسط «متحدان» خود تحویل دربار شاهنشاه شدند.%
  \footnote{\lr{Briant, 2002, pp.~656--675.}}
  \item \textbf{قدرت بازسازی مرکز:} اردشیر دوم با ترکیبی از دیپلماسی و زور نظامی شورش را فرونشاند.
\end{enumerate}
\end{tcolorbox}

\subsection{اسکندر و مسئله‌ی «دعوت‌شدگان»}
\label{subsec:ch04-alexander}

حمله‌ی \histref{اسکندر مقدونی}{۳۳۴–۳۲۳ ق.م} به شاهنشاهی هخامنشی، یکی از پیچیده‌ترین موارد تاریخی از منظر نجات‌طلبی است. پرسش کلیدی آن است: آیا بخشی از نخبگان یا مردم ایران از اسکندر «دعوت» کردند یا استقبال نمودند؟

\subsubsection{شواهد له}

\begin{enumerate}[nosep, label=\textbf{\arabic*.}]
  \item \textbf{تسلیم بدون مقاومت برخی شهرها:} \thinker{آریان} (\lr{Arrian}) گزارش می‌دهد که بسیاری از شهرهای آسیای صغیر و بابل بدون مقاومت تسلیم شدند و حتی اسکندر را «آزادکننده» خواندند.%
  \footnote{\lr{Arrian. \textit{Anabasis Alexandri}. III.16. Trans. P.A. Brunt. Loeb Classical Library, 1976.}}
  \item \textbf{همکاری ساتراپ‌ها:} برخی ساتراپ‌ها مانند \enemy{مازائوس} (ساتراپ بابل) بدون جنگ تسلیم شدند و حتی در دستگاه اداری اسکندر مقام یافتند.%
  \footnote{\lr{Bosworth, A.B. \textit{Conquest and Empire: The Reign of Alexander the Great}. Cambridge UP, 1988, pp.~86--90.}}
  \item \textbf{نارضایتی از داریوش سوم:} شواهدی وجود دارد که برخی نخبگان ایرانی از ضعف نظامی و مدیریتی \enemy{داریوش سوم} ناراضی بودند — قتل او توسط \enemy{بسوس}، ساتراپ باختر، نشانه‌ی این نارضایتی است.%
  \footnote{\lr{Briant, 2002, pp.~817--869.}}
\end{enumerate}

\subsubsection{شواهد علیه}

\begin{enumerate}[nosep, label=\textbf{\arabic*.}]
  \item \textbf{مقاومت‌های سرسختانه:} نبردهای ایسوس و گوگمل نشان‌دهنده‌ی مقاومت جدی ارتش هخامنشی بود.
  \item \textbf{مقاومت محلی:} برخی مناطق — به‌ویژه سغد و باختر — سال‌ها در برابر اسکندر مقاومت کردند. \histref{اسپیتامنِس}{۳۲۹--۳۲۸ ق.م} رهبر شورش سغدیان، مقاومتی طولانی و خونین ترتیب داد.%
  \footnote{\lr{Holt, Frank L. \textit{Into the Land of Bones: Alexander the Great in Afghanistan}. U of California P, 2005.}}
  \item \textbf{ماهیت «تسلیم»:} \thinker{پی‌یر بریان} استدلال می‌کند که تسلیم بسیاری از شهرها نه «دعوت» بلکه «تمکین عملی» (\lr{pragmatic submission}) در برابر قدرت برتر نظامی بود.%
  \footnote{\lr{Briant, 2002, pp.~840--855.}}
\end{enumerate}

\begin{figure}[htbp]
\centering
\begin{tikzpicture}[
  every node/.style={font=\scriptsize},
  node distance=0.5cm,
]

% نقشه‌ی شماتیک مسیر اسکندر
\draw[line width=1.5pt, PrimaryMid, ->]
  (0,0) node[left, font=\small\bfseries] {\rl{مقدونیه}} --
  (2,0.5) node[above] {\rl{هلسپونت}} --
  (4,0.8) node[above, text=AccentGreen] {\rl{تسلیم: ساردیس}} --
  (5.5,0.3) node[below, text=AccentRed] {\rl{نبرد ایسوس}} --
  (7,-0.5) node[below, text=AccentGreen] {\rl{تسلیم: مصر}} --
  (9,0.2) node[above, text=AccentRed] {\rl{نبرد گوگمل}} --
  (11,0) node[above, text=AccentGreen] {\rl{تسلیم: بابل}} --
  (12.5,0.5) node[above, text=AccentGreen] {\rl{شوش، تخت‌جمشید}} --
  (14,1) node[right, text=AccentRed] {\rl{مقاومت: سغد و باختر}};

% راهنما
\node[text=AccentGreen, font=\scriptsize\bfseries] at (3,-1.5) {\rl{سبز = تسلیم/استقبال}};
\node[text=AccentRed, font=\scriptsize\bfseries] at (9,-1.5) {\rl{قرمز = مقاومت/نبرد}};

\end{tikzpicture}
\caption{نقشه‌ی شماتیک مسیر اسکندر: مقاومت و تسلیم}
\label{fig:ch04-alexander-route}
\end{figure}

\begin{tcolorbox}[casebox, title={تحلیل شش‌وجهی: حمله‌ی اسکندر و واکنش‌های داخلی}]
\begin{itemize}[nosep, rightmargin=1em]
  \item \textbf{ساختاری:} شاهنشاهی هخامنشی در اواخر عمر خود دچار ضعف ساختاری شده بود — جنگ‌های داخلی، شورش‌های مکرر ساتراپ‌ها و بحران جانشینی. اما این ضعف به‌معنای «دعوت» از اسکندر نبود.
  \item \textbf{حقوقی-هنجاری:} مفهوم مدرن حاکمیت ملی وجود نداشت. تسلیم ساتراپ‌ها بیشتر ناشی از محاسبه‌ی عملی بود تا حقوقی.
  \item \textbf{اقتصادی:} اسکندر ثروت هنگفت خزانه‌های هخامنشی (تخت‌جمشید، شوش، اکباتان) را غارت کرد — حدود ۱۸۰٬۰۰۰ تالنت نقره.%
  \footnote{\lr{Hammond, N.G.L. \textit{The Genius of Alexander the Great}. U of North Carolina P, 1997, p.~142.}}
  \item \textbf{نظامی:} برتری تاکتیکی مقدونی‌ها قاطع بود. مقاومت مؤثر فقط در مناطق کوهستانی و دورافتاده ممکن شد.
  \item \textbf{روان‌شناختی:} «شوک و بهت» (\lr{shock and awe}) حاصل از سرعت پیشروی اسکندر، مقاومت روانی بسیاری را درهم شکست.
  \item \textbf{ژئوپولیتیک:} حمله‌ی اسکندر در بستر رقابت دیرینه‌ی یونان-ایران و با ایدئولوژی «انتقام» از حمله‌ی خشایارشا (۴۸۰ ق.م) صورت گرفت.%
  \footnote{\lr{Flower, Michael A. "Alexander the Great and Panhellenism." In \textit{Alexander the Great in Fact and Fiction}, ed. Bosworth and Baynham. Oxford UP, 2000, pp.~96--135.}}
  \item \textbf{نتیجه (طیف):} \textbf{E–F}. فروپاشی کامل شاهنشاهی، غارت ثروت ملی، و تحمیل فرهنگ یونانی (هلنیسم) — هرچند بخشی از ساختار اداری محلی حفظ شد.
\end{itemize}
\end{tcolorbox}

% ─────────────────────────────────────────────────────────────────────────
\section{دوره‌ی اشکانی (پارتی): ۲۴۷ ق.م — ۲۲۴ م.}
\label{sec:ch04-parthian}
% ─────────────────────────────────────────────────────────────────────────

\subsection{ساختار فدرالی و آسیب‌پذیری‌ها}
\label{subsec:ch04-parthian-structure}

شاهنشاهی اشکانی ساختاری \concept{فدرالی} و \concept{فئودالی} داشت که در آن خاندان‌های بزرگ (سورن، کارن، مهران و غیره) قدرت محلی قابل‌توجهی داشتند. \thinker{ملکم کالج} (\lr{Malcolm Colledge}) این ساختار را «کنفدراسیون پادشاهی‌ها» خوانده است.%
\footnote{\lr{Colledge, Malcolm. \textit{The Parthians}. Thames and Hudson, 1967, p.~37.}}

این ساختار، امکان نجات‌طلبی را افزایش می‌داد زیرا:
\begin{itemize}[nosep, rightmargin=1em]
  \item خاندان‌های ناراضی می‌توانستند مستقلاً با قدرت خارجی (روم) وارد مذاکره شوند.
  \item مدعیان تاج‌وتخت اشکانی اغلب از روم درخواست کمک می‌کردند.
  \item ارمنستان به‌عنوان دولت حائل، محل تلاقی نجات‌طلبی‌های ایرانی و رومی بود.
\end{itemize}

\subsection{نمونه‌ی کلیدی: تیرداد و روم}
\label{subsec:ch04-tiridates}

در سال‌های ۳۶–۲۶ ق.م، \histref{تیرداد دوم}{شاهزاده‌ی اشکانی} که از برادرش \enemy{فرهاد چهارم} شکست خورده بود، به دربار \histref{اوکتاویان (آگوستوس)}{۳۶ ق.م} در روم پناه برد و از او خواست تا نظامی به ایران بفرستد و او را بر تخت بنشاند. آگوستوس با احتیاط عمل کرد و از مداخله‌ی مستقیم نظامی خودداری نمود، اما از تیرداد به‌عنوان ابزار فشار دیپلماتیک علیه فرهاد استفاده کرد.%
\footnote{\lr{Dąbrowa, Edward. "Roman Policy in Transcaucasia from Pompey to Domitian." In \textit{The Roman and Byzantine Army in the East}, ed. Dąbrowa. Jagiellonian UP, 1994, pp.~67--76.}}

\begin{tcolorbox}[defbox, title={درس تاریخی}]
مورد تیرداد نشان می‌دهد که حتی ابرقدرت‌های باستانی نیز در مداخله‌ی نظامی مستقیم محتاط بودند. آگوستوس دریافت که حفظ تیرداد به‌عنوان «تهدید بالقوه» سودمندتر از جنگ مستقیم با پارت بود. این الگو — \concept{استفاده‌ی ابزاری از مدعیان تبعیدی بدون مداخله‌ی مستقیم} — در سراسر تاریخ تکرار شده است.
\end{tcolorbox}

\subsection{جدول تطبیقی: مدعیان تبعیدی اشکانی در روم}
\label{subsec:ch04-parthian-exiles}

\begin{longtable}{
  >{\raggedleft\arraybackslash\bfseries}p{2.5cm}
  >{\raggedleft\arraybackslash}p{2cm}
  >{\raggedleft\arraybackslash}p{3cm}
  >{\raggedleft\arraybackslash}p{2.5cm}
  >{\raggedleft\arraybackslash}p{3cm}}
\caption{مدعیان اشکانی در روم و سرنوشت آنها}\label{tab:ch04-parthian-exiles}\\
\toprule
\rowcolor{PrimaryDeep}
\textcolor{white}{\textbf{مدعی}} &
\textcolor{white}{\textbf{تاریخ}} &
\textcolor{white}{\textbf{حامی رومی}} &
\textcolor{white}{\textbf{مداخله‌ی نظامی}} &
\textcolor{white}{\textbf{نتیجه}} \\
\midrule
\endfirsthead
\toprule
\rowcolor{PrimaryDeep}
\textcolor{white}{\textbf{مدعی}} &
\textcolor{white}{\textbf{تاریخ}} &
\textcolor{white}{\textbf{حامی رومی}} &
\textcolor{white}{\textbf{مداخله‌ی نظامی}} &
\textcolor{white}{\textbf{نتیجه}} \\
\midrule
\endhead
\bottomrule
\endlastfoot

تیرداد دوم &
ح. ۳۰ ق.م &
آگوستوس &
خیر (ابزار فشار) &
شکست؛ بازگشت نکرد%
\footnote{\lr{Justin. \textit{Epitome of Trogus}. XLII.5.}} \\
\rowcolor{NeutralLight}
ونونس اول &
ح. ۱۲ م. &
آگوستوس/تیبریوس &
بله (نیمه‌مستقیم) &
تاج‌گذاری اما سقوط سریع%
\footnote{\lr{Tacitus. \textit{Annals}. II.1--4. Trans. A.J. Woodman. Hackett, 2004.}} \\
مهرداد &
ح. ۳۵–۳۶ م. &
تیبریوس &
خیر &
شکست \\
\rowcolor{NeutralLight}
تیرداد سوم ارمنستان &
ح. ۶۳ م. &
نرون &
مذاکره (صلح راندیا) &
تاج‌گذاری نمادین در روم%
\footnote{\lr{Dąbrowa, Edward. "Gaius Cassius Longinus's Command in Arsacid Context." \textit{Electrum} 19 (2012): 157--165.}} \\
خسرو (اسوارس) &
ح. ۱۰۹ م. &
تراژان &
بله (لشکرکشی بزرگ) &
فتح موقت بین‌النهرین؛ عقب‌نشینی%
\footnote{\lr{Bennett, Julian. \textit{Trajan: Optimus Princeps}. Routledge, 1997, pp.~183--204.}} \\

\end{longtable}

% ─────────────────────────────────────────────────────────────────────────
\section{دوره‌ی ساسانی (۲۲۴–۶۵۱ م.)}
\label{sec:ch04-sasanian}
% ─────────────────────────────────────────────────────────────────────────

\subsection{ساختار قدرت و خطوط شکاف}
\label{subsec:ch04-sasanian-structure}

شاهنشاهی ساسانی نسبت به اشکانی متمرکزتر بود اما خطوط شکاف مهمی داشت:

\begin{itemize}[nosep, rightmargin=1em]
  \item \textbf{شکاف مذهبی:} زرتشتیان ارتدوکس در مقابل مسیحیان (به‌ویژه در بین‌النهرین)، مانویان و مزدکیان.
  \item \textbf{شکاف قومی:} ارمنی‌ها، آرامی‌ها، عرب‌های مرزنشین (لخمیان) و کردها.
  \item \textbf{شکاف طبقاتی:} نجبا و موبدان در مقابل دهقانان و رعایا.
\end{itemize}

\thinker{تورج دریایی} (\lr{Touraj Daryaee}) تأکید می‌کند که «ساسانیان با وجود تمرکزگرایی ظاهری، همواره با چالش نخبگان محلی و اقلیت‌های مذهبی روبه‌رو بودند.»%
\footnote{\lr{Daryaee, Touraj. \textit{Sasanian Persia: The Rise and Fall of an Empire}. I.B. Tauris, 2009, pp.~28--45.}}

\subsection{مسیحیان بین‌النهرین و استمداد از روم-بیزانس}
\label{subsec:ch04-christians}

جمعیت بزرگ مسیحی بین‌النهرین — به‌ویژه در شهرهای نصیبین، اِدِسا و سلوکیه-تیسفون — نقش مهمی در تنش‌های ساسانی-رومی داشت. پس از مسیحی‌شدن امپراتوری روم (۳۱۳ م.)، مسیحیان ایران به‌طور فزاینده‌ای مظنون به وفاداری به روم شدند.%
\footnote{\lr{Brock, Sebastian. "Christians in the Sasanian Empire." \textit{Studies in Church History} 18 (1982): 1--19.}}

\concept{پادافره‌ی بزرگ} (\lr{Great Persecution}) شاپور دوم علیه مسیحیان (ح. ۳۴۰–۳۷۹ م.) بخشاً واکنش به این مسئله بود. شاپور دوم در نامه‌ای به \thinker{شمعون بار صبّاعه} (کاتولیکوس مسیحیان ایران) نوشت: «شمعون دوست قیصر ماست و علیه کشور ما جاسوسی می‌کند.»%
\footnote{\lr{Asmussen, J.P. "Christians in Iran." In \textit{Cambridge History of Iran}, vol.~3(2), ed. Yarshater. Cambridge UP, 1983, pp.~924--948.}}

\begin{tcolorbox}[casebox, title={تحلیل: مسیحیان ساسانی — نجات‌طلبی یا تهمت؟}]
آیا مسیحیان ایران واقعاً از روم «دعوت» می‌کردند؟ شواهد مبهم‌اند:
\begin{itemize}[nosep, rightmargin=1em]
  \item \textbf{شواهد له:} نامه‌نگاری‌های برخی اسقف‌ها با دربار بیزانس؛ فرار مسیحیان ایرانی به قلمرو روم.
  \item \textbf{شواهد علیه:} کلیسای شرق (\lr{Church of the East}) در شورای مرو (۴۲۴ م.) رسماً استقلال خود از کلیسای رومی را اعلام کرد — دقیقاً برای رفع اتهام وفاداری به روم.%
  \footnote{\lr{Baum, Wilhelm, and Dietmar W. Winkler. \textit{The Church of the East}. RoutledgeCurzon, 2003, pp.~19--22.}}
  \item \textbf{نکته‌ی کلیدی:} حتی اگر «دعوت» صریحی وجود نداشت، \concept{ادراک حکومت از نجات‌طلبی} (\lr{perceived salvation-seeking}) به اندازه‌ی واقعیت آن مهم بود و پیامدهای خونینی داشت.
\end{itemize}
\end{tcolorbox}

\subsection{لخمیان: عرب‌های مرزنشین و دوگانگی وفاداری}
\label{subsec:ch04-lakhmids}

\concept{لخمیان} حکومت عرب دست‌نشانده‌ی ساسانی در حیره (عراق امروزی) بودند که نقش سپر مرزی علیه قبایل عرب صحرانشین و غسانیان (متحد بیزانس) را ایفا می‌کردند. در سال ۶۰۲ م.، \enemy{خسرو دوم} آخرین پادشاه لخمی (\enemy{نعمان سوم}) را بر‌کنار و اعدام کرد.%
\footnote{\lr{Shahîd, Irfan. \textit{Byzantium and the Arabs in the Sixth Century}. Dumbarton Oaks, 1995, vol.~I, pp.~305--328.}}

این تصمیم فاجعه‌بار بود: \textbf{قبایل عرب بدون واسطه‌ی لخمیان، دیگر دلیلی برای وفاداری به ساسانیان نداشتند.} در نبرد ذوقار (ح. ۶۰۴–۶۱۰ م.) قبایل عرب برای نخستین بار ارتش ساسانی را شکست دادند — رویدادی که بعدها مسلمانان آن را پیش‌درآمد فتوحات خود دانستند.%
\footnote{طبری، محمد بن جریر. \textit{تاریخ الرسل و الملوک}. تحقیق محمد ابوالفضل ابراهیم. دار المعارف، ۱۹۶۷، ج ۲، ص ۱۹۵--۲۰۲.}

\begin{tcolorbox}[enemybox, title={درس تاریخی: حذف واسطه‌ها و پیامدهای ناخواسته}]
حذف لخمیان توسط خسرو دوم نمونه‌ای است از آنچه در علوم سیاسی \concept{پیامد ناخواسته} (\lr{Unintended Consequences}) نامیده می‌شود: عملی که قرار بود قدرت مرکزی را تحکیم کند، عملاً مرز جنوبی شاهنشاهی را بی‌دفاع گذاشت و زمینه‌ساز فتح عرب شد.%
\footnote{\lr{Howard-Johnston, James. \textit{Witnesses to a World Crisis: Historians and Histories of the Middle East in the Seventh Century}. Oxford UP, 2010, pp.~436--445.}}
\end{tcolorbox}

\subsection{فروپاشی ساسانی و مسئله‌ی «همکاری» با اعراب}
\label{subsec:ch04-arab-conquest}

\subsubsection{آیا بخشی از ایرانیان از اعراب «استقبال» کردند؟}

این پرسشی حساس و پربحث است. شواهد تاریخی نشان‌دهنده‌ی طیفی از واکنش‌هاست:

\begin{longtable}{
  >{\raggedleft\arraybackslash\bfseries}p{2.5cm}
  >{\raggedleft\arraybackslash}p{4cm}
  >{\raggedleft\arraybackslash}p{3cm}
  >{\raggedleft\arraybackslash}p{3cm}}
\caption{واکنش‌های محلی به فتح عرب در ایران}\label{tab:ch04-arab-reactions}\\
\toprule
\rowcolor{PrimaryDeep}
\textcolor{white}{\textbf{منطقه}} &
\textcolor{white}{\textbf{واکنش}} &
\textcolor{white}{\textbf{عامل اصلی}} &
\textcolor{white}{\textbf{منبع}} \\
\midrule
\endfirsthead
\toprule
\rowcolor{PrimaryDeep}
\textcolor{white}{\textbf{منطقه}} &
\textcolor{white}{\textbf{واکنش}} &
\textcolor{white}{\textbf{عامل اصلی}} &
\textcolor{white}{\textbf{منبع}} \\
\midrule
\endhead
\bottomrule
\endlastfoot

بین‌النهرین (سواد) &
تسلیم سریع؛ دهقانان برخی شهرها دروازه‌ها را گشودند &
نارضایتی از مالیات‌های سنگین؛ جمعیت بزرگ مسیحی و یهودی &
طبری، بلاذری%
\footnote{بلاذری، احمد بن یحیی. \textit{فتوح البلدان}. تحقیق عبدالله انیس الطباع. مؤسسه المعارف، ۱۹۸۷، ص ۲۵۹--۲۸۰.} \\
\rowcolor{NeutralLight}
فارس &
مقاومت سرسختانه‌ی نظامی (نبرد نهاوند ۶۴۲) اما سپس تسلیم &
وفاداری به ساختار ساسانی &
طبری \\
خراسان &
مقاومت طولانی و مکرر؛ شورش‌های پیاپی تا اواخر سده‌ی هفتم &
هویت ایرانی قوی؛ فاصله از مرکز خلافت &
بلاذری، ابن‌اثیر%
\footnote{ابن‌اثیر، عزالدین. \textit{الکامل فی التاریخ}. دار صادر، ۱۹۶۵، ج ۲، ص ۳۴۵--۳۸۰.} \\
\rowcolor{NeutralLight}
طبرستان و گیلان &
مقاومت مستمر تا ح. ۱۴۰ هجری؛ حکومت‌های محلی مستقل &
جغرافیای کوهستانی؛ استقلال‌طلبی دیرینه &
ابن‌اسفندیار%
\footnote{ابن‌اسفندیار، محمد بن حسن. \textit{تاریخ طبرستان}. تصحیح عباس اقبال. کلاله‌ی خاور، ۱۳۲۰.} \\
سیستان &
مقاومت و سپس تسلیم مشروط (صلح‌نامه) &
سنت استقلال محلی &
تاریخ سیستان%
\footnote{ناشناس. \textit{تاریخ سیستان}. تصحیح محمدتقی بهار. کلاله‌ی خاور، ۱۳۱۴، ص ۸۰--۹۵.} \\

\end{longtable}

\subsubsection{تحلیل: «دعوت» یا «تمکین»؟}

\thinker{پاتریشیا کرون} (\lr{Patricia Crone}) و \thinker{مارتین هیندز} (\lr{Martin Hinds}) در پژوهش تأثیرگذار خود استدلال می‌کنند که تسلیم بسیاری از شهرهای ایرانی نه «استقبال از اسلام» بلکه نتیجه‌ی \concept{فروپاشی ساختار فرماندهی ساسانی} پس از مرگ یزدگرد سوم بود.%
\footnote{\lr{Crone, Patricia, and Martin Hinds. \textit{God's Caliph: Religious Authority in the First Centuries of Islam}. Cambridge UP, 1986, pp.~24--42.}}

\thinker{پروانه پورشریعتی} (\lr{Parvaneh Pourshariati}) در اثر مهم «زوال و فروپاشی شاهنشاهی ساسانی» نظریه‌ی نوینی ارائه می‌دهد: فروپاشی ساسانی نه نتیجه‌ی حمله‌ی عرب بلکه نتیجه‌ی \concept{جنگ داخلی میان خاندان‌های پارتی و ساسانی} بود. برخی خاندان‌های پارتی — به‌ویژه خاندان مهران — عملاً با اعراب «همکاری» کردند نه از سر نجات‌طلبی، بلکه برای تسویه‌حساب‌های داخلی.%
\footnote{\lr{Pourshariati, Parvaneh. \textit{Decline and Fall of the Sasanian Empire}. I.B. Tauris, 2008, pp.~175--230.}}

\begin{tcolorbox}[casebox, title={تحلیل شش‌وجهی: فتح عرب ایران (۶۳۶–۶۵۱)}]
\begin{itemize}[nosep, rightmargin=1em]
  \item \textbf{ساختاری:} بحران عمیق جانشینی (۱۲ شاه در ۴ سال: ۶۲۸–۶۳۲)؛ جنگ داخلی؛ خستگی از جنگ‌های طولانی با بیزانس.%
  \footnote{\lr{Howard-Johnston, 2010, pp.~436--460.}}
  \item \textbf{حقوقی-هنجاری:} مفهوم «جهاد» مبنای ایدئولوژیک فتح بود — نه «دعوت» ایرانیان. اما صلح‌نامه‌ها (\lr{sulh}) با برخی شهرها نشان‌دهنده‌ی عنصر رضایت محلی است.
  \item \textbf{اقتصادی:} جزیه و خراج — بار مالی سنگین بر غیرمسلمانان. اما برای دهقانان ایرانی، مالیات‌های عرب در ابتدا سبک‌تر از مالیات‌های ساسانی بود.%
  \footnote{دینوری، ابوحنیفه. \textit{اخبار الطوال}. تحقیق عبدالمنعم عامر. وزارة الثقافة، ۱۹۶۰، ص ۱۳۵--۱۴۲.}
  \item \textbf{نظامی:} شکست قاطع در نبردهای قادسیه (۶۳۶) و نهاوند (۶۴۲)؛ فقدان فرماندهی متمرکز پس از مرگ رستم فرخزاد.
  \item \textbf{روان‌شناختی:} «شوک تمدنی» — فروپاشی جهان‌بینی ساسانی که شاهنشاه را نماینده‌ی اهورامزدا می‌دانست.%
  \footnote{\lr{Daryaee, 2009, pp.~112--128.}}
  \item \textbf{ژئوپولیتیک:} ضعف هم‌زمان ساسانیان و بیزانس پس از جنگ‌های فرسایشی ۶۰۲–۶۲۸، خلأ قدرت منطقه‌ای ایجاد کرده بود.
  \item \textbf{نتیجه (طیف):} \textbf{F}. فروپاشی کامل شاهنشاهی، تغییر دین و فرهنگ، از دست‌رفتن استقلال سیاسی برای قرن‌ها. البته در بلندمدت (بیش از دو قرن)، هویت ایرانی بازسازی شد — اما در قالبی بنیاداً متفاوت.
\end{itemize}
\end{tcolorbox}

% ─────────────────────────────────────────────────────────────────────────
\section{خط زمانی تطبیقی: نجات‌طلبی در ایران باستان}
\label{sec:ch04-timeline}
% ─────────────────────────────────────────────────────────────────────────

\begin{figure}[htbp]
\centering
\begin{tikzpicture}[
  x=0.02\linewidth,
  every node/.style={font=\scriptsize},
]

% خط اصلی
\draw[line width=2.5pt, PrimaryMid] (0,0) -- (120,0);

% دوره‌ها
\fill[EraAncient, opacity=0.2] (0,-0.6) rectangle (30,0.6);
\fill[EraMedieval!70, opacity=0.2] (30,-0.6) rectangle (70,0.6);
\fill[EraEarlyMod!70, opacity=0.2] (70,-0.6) rectangle (120,0.6);

% برچسب دوره‌ها
\node[era label, fill=EraAncient] at (15,1.5) {\rl{هخامنشی}};
\node[era label, fill=EraMedieval!80] at (50,1.5) {\rl{اشکانی}};
\node[era label, fill=EraEarlyMod!80] at (95,1.5) {\rl{ساسانی}};

% رویدادها - بالا
\node[event node, above=2cm] (e1) at (5,0)
  {\rl{فتح بابل توسط کورش}\\%
   \rl{۵۳۹ ق.م}};
\draw[thin, PrimaryLight] (5,0) -- (e1);

\node[event node, above=2cm] (e2) at (22,0)
  {\rl{شورش ساتراپ‌ها}\\%
   \rl{۳۶۶--۳۶۰ ق.م}};
\draw[thin, PrimaryLight] (22,0) -- (e2);

\node[event node, above=2cm] (e3) at (38,0)
  {\rl{تیرداد و آگوستوس}\\%
   \rl{ح. ۳۰ ق.م}};
\draw[thin, PrimaryLight] (38,0) -- (e3);

\node[event node, above=2cm] (e4) at (60,0)
  {\rl{ونونس و تیبریوس}\\%
   \rl{ح. ۱۲ م.}};
\draw[thin, PrimaryLight] (60,0) -- (e4);

% رویدادها - پایین
\node[event node, below=2cm] (e5) at (28,0)
  {\rl{حمله‌ی اسکندر}\\%
   \rl{۳۳۴--۳۲۳ ق.م}};
\draw[thin, AccentRed] (28,0) -- (e5);

\node[event node, below=2cm] (e6) at (80,0)
  {\rl{پادافره‌ی مسیحیان}\\%
   \rl{ح. ۳۴۰ م.}};
\draw[thin, AccentRed] (80,0) -- (e6);

\node[event node, below=2cm] (e7) at (100,0)
  {\rl{حذف لخمیان}\\%
   \rl{۶۰۲ م.}};
\draw[thin, AccentRed] (100,0) -- (e7);

\node[event node, below=2cm] (e8) at (115,0)
  {\rl{فتح عرب}\\%
   \rl{۶۳۶--۶۵۱ م.}};
\draw[thin, AccentRed] (115,0) -- (e8);

\end{tikzpicture}
\caption{خط زمانی نجات‌طلبی و مداخلات خارجی در ایران باستان}
\label{fig:ch04-timeline}
\end{figure}

% ─────────────────────────────────────────────────────────────────────────
\section{الگوهای تکرارشونده در ایران باستان}
\label{sec:ch04-patterns}
% ─────────────────────────────────────────────────────────────────────────

از بررسی تاریخ ایران باستان، چندین الگوی تکرارشونده قابل استخراج است:

\begin{tcolorbox}[wavebox, title={الگوهای استخراج‌شده}]
\begin{enumerate}[nosep, label=\textbf{الگوی \arabic*:}]
  \item \textbf{ساتراپ‌های ناراضی به‌عنوان عاملان اصلی نجات‌طلبی:} در هر سه دوره (هخامنشی، اشکانی، ساسانی)، نخبگان محلی و مرزنشین بیشتر از توده‌ی مردم به بیگانه متوسل شده‌اند.

  \item \textbf{اقلیت‌های مذهبی به‌عنوان «پل ارتباطی»:} مسیحیان بین‌النهرین، یهودیان بابل و بعدها مانویان، به‌دلیل پیوندهای فرامرزی‌شان، مظنون به نجات‌طلبی بودند — خواه واقعی خواه فرضی.

  \item \textbf{استفاده‌ی ابزاری ابرقدرت‌ها از مدعیان تبعیدی:} روم از شاهزادگان اشکانی و ساسانی به‌عنوان ابزار فشار استفاده می‌کرد بدون آنکه لزوماً به تعهداتش عمل کند.

  \item \textbf{حذف واسطه‌ها با پیامدهای فاجعه‌بار:} حذف لخمیان توسط خسرو دوم مرز جنوبی را بی‌دفاع کرد — الگویی که در تاریخ ایران تکرار خواهد شد (مثلاً حذف خوانین مرزی توسط رضاشاه).

  \item \textbf{تفاوت «تمکین عملی» و «دعوت»:} بسیاری از مواردی که در منابع به‌عنوان «استقبال» گزارش شده‌اند، در واقع تسلیم واقع‌بینانه در برابر قدرت برتر نظامی بوده‌اند و نباید با «نجات‌طلبی» خلط شوند.
\end{enumerate}
\end{tcolorbox}

% ─────────────────────────────────────────────────────────────────────────
\section{جدول تطبیقی نهایی: نجات‌طلبی در سه دوره‌ی ایران باستان}
\label{sec:ch04-final-table}
% ─────────────────────────────────────────────────────────────────────────

\begin{landscape}
\begin{table}[htbp]
\centering
\caption{جدول تطبیقی نجات‌طلبی در ایران باستان}\label{tab:ch04-comparative}
\begin{adjustbox}{max width=\linewidth}
\begin{tabular}{
  >{\raggedleft\arraybackslash\bfseries\small}p{2.2cm}
  >{\raggedleft\arraybackslash\small}p{2.5cm}
  >{\raggedleft\arraybackslash\small}p{2cm}
  >{\raggedleft\arraybackslash\small}p{2cm}
  >{\raggedleft\arraybackslash\small}p{2cm}
  >{\raggedleft\arraybackslash\small}p{2.5cm}
  >{\raggedleft\arraybackslash\small}p{2cm}
  >{\raggedleft\arraybackslash\small}p{2cm}}
\toprule
\rowcolor{PrimaryDeep}
\textcolor{white}{مورد} &
\textcolor{white}{دعوت‌کننده} &
\textcolor{white}{مداخله‌گر} &
\textcolor{white}{نوع مداخله} &
\textcolor{white}{مبنای مشروعیت} &
\textcolor{white}{نتیجه‌ی فوری} &
\textcolor{white}{نتیجه‌ی بلندمدت} &
\textcolor{white}{جایگاه طیفی} \\
\midrule

کورش-بابل &
کاهنان بابلی &
ایران (هخامنشی) &
نظامی مستقیم &
حق الهی (مردوک) &
سقوط نبونید &
ادغام در شاهنشاهی &
A–B \\
\rowcolor{NeutralLight}
شورش ساتراپ‌ها &
ساتراپ‌های آسیای صغیر &
اسپارت، مصر &
نظامی غیرمستقیم &
منافع متقابل &
شکست شورش &
تحکیم مرکز &
شکست \\
اسکندر &
تسلیم برخی ساتراپ‌ها &
مقدونیه &
نظامی مستقیم &
«انتقام» یونانی &
فروپاشی هخامنشی &
هلنیسم، از دست‌رفتن استقلال &
E–F \\
\rowcolor{NeutralLight}
مدعیان اشکانی &
شاهزادگان تبعیدی &
روم &
دیپلماتیک + نظامی &
ابزار فشار &
نتایج مختلط &
بی‌ثباتی مزمن &
C–D \\
مسیحیان ساسانی &
(مظنون) اسقف‌ها &
بیزانس &
فرهنگی-مذهبی &
هم‌کیشی &
سرکوب مسیحیان &
استقلال کلیسای شرق &
شکست (برای مسیحیان) \\
\rowcolor{NeutralLight}
فتح عرب &
(نه دعوت صریح) &
خلافت راشدین &
نظامی مستقیم &
جهاد &
فروپاشی ساسانی &
تغییر تمدنی &
F \\

\bottomrule
\end{tabular}
\end{adjustbox}
\end{table}
\end{landscape}

% ─────────────────────────────────────────────────────────────────────────
\section{جمع‌بندی فصل}
\label{sec:ch04-summary}
% ─────────────────────────────────────────────────────────────────────────

بررسی تاریخ ایران باستان نشان داد که:

\begin{enumerate}[nosep, label=\textbf{\arabic*.}]
  \item نجات‌طلبی صریح — یعنی «دعوت» آگاهانه و سازمان‌یافته از نیروی خارجی — در دوران باستان نادر بود. آنچه بیشتر دیده می‌شود، \concept{تمکین عملی} در برابر قدرت برتر یا \concept{استفاده‌ی ابزاری} ابرقدرت‌ها از مدعیان داخلی است.

  \item ساختار فدرالی/فئودالی امپراتوری‌ها، نخبگان محلی و مرزنشین را به بازیگران اصلی نجات‌طلبی تبدیل می‌کرد — نه توده‌ی مردم.

  \item اقلیت‌های مذهبی حتی بدون نجات‌طلبی واقعی، مظنون به آن بودند و بهای سنگینی می‌پرداختند.

  \item در هیچ‌یک از موارد بررسی‌شده، مداخله‌ی خارجی به نفع بلندمدت دعوت‌کنندگان (اگر وجود داشتند) تمام نشد — حتی ورود کورش به بابل که «موفق‌ترین» نمونه بود، به از دست‌رفتن استقلال سیاسی بابل انجامید.

  \item الگوی «مست گیرند، در شهر هر آنکه هست گیرند» در تمامی موارد صادق بود: مداخله‌گر پس از ورود، منافع خود را بر خواسته‌ی دعوت‌کنندگان ترجیح داد.
\end{enumerate}

\cleardoublepage